% This is the aspauthor.tex LaTeX file
% Copyright 2021, Astronomical Society of the Pacific Conference Series
% Revision:  29 November 2021

% To compile, at the command line positioned at this folder, type:
% pdflatex aspauthor
% bibtex aspauthor
% pdflatex aspauthor
% pdflatex aspauthor
% This will create a file called aspauthor.pdf.

\documentclass[11pt,twoside]{article}
\usepackage{asp2023}

\aspSuppressVolSlug
\resetcounters

\bibliographystyle{asp2021}

\markboth{Gangat and Schwartz}{Building Confidence in PAWS with Playwright}

\begin{document}

\title{Building Confidence in the MeerKAT Proposal Submission System with Playwright}
\author{A.~Gangat and R.~L.~Schwartz}
\affil{South African Radio Astronomy Observatory, Cape Town, South Africa; \email{agangat@sarao.ac.za}; \email{rschwartz@sarao.ac.za}}


% This section is for ADS Processing.  There must be one line per author.
\paperauthor{A.~Gangat}{agangat@sarao.ac.za}{}{South African Radio Astronomy Observatory}{Software}{Cape Town}{}{}{South Africa}
\paperauthor{R.~L.~Schwartz}{rschwartz@sarao.ac.za}{}{South African Radio Astronomy Observatory}{Software}{Cape Town}{}{}{South Africa}

% Author-index entries required by the ADASS template.
%\aindex{Gangat,~A.}
%\aindex{Schwartz,~R.~L.}

\begin{abstract}
The Proposal Workflow System (PAWS) manages proposal submissions for the MeerKAT Radio Telescope, making reliability an important part of the user experience. This work introduces automated regression testing with Playwright to check that existing functionality is maintained as the application evolves.

Tests create draft proposals and exercise key stages of the submission process. They also cover common form field validation scenarios and behaviours. Reusable helpers and test documents support repeatability and cleanup.

The suite provides repeatable evidence that selected behaviours remain intact after application changes, including AI-assisted changes. It builds confidence in the tested workflows while complementing human review and other testing approaches, and provides a foundation for broader coverage.
\end{abstract}

\section{Introduction and Motivation}

Preparing a proposal for observing time on the MeerKAT Radio Telescope requires team details, science information, observing requirements, data-management arrangements and supporting documents. PAWS collects this information for technical, science and panel reviews \citep{SchwartzBaloyiSithole2024PAWSSubmission,SchwartzBaloyi2024PAWSReview}, making reliable submission workflows essential.\footnote{MeerKAT: \url{https://www.sarao.ac.za/science/meerkat/}.}

PAWS replaced a largely manual submission process. Tight timelines initially limited unit testing, with features verified manually; later calls required a more customisable, in-house form \citep{SchwartzBaloyiSithole2024PAWSSubmission}. Between 2020 and 2023, submitted proposals increased from 112 to 170 and requested telescope time from 3343 to 7554 hours \citep{SchwartzBaloyiSithole2024PAWSSubmission}. By 2024, PAWS supported 71 science reviewers and 772 reviews, while maintaining legacy data increased code complexity \citep{SchwartzBaloyi2024PAWSReview}.

During each call for proposals, requests from the Chief Scientist continue to refine submission workflows, while earlier proposal data must remain accessible for reuse and statistics. Changes to requirements, validation and frontend components, including Angular Material's migration to Material Design Components (MDC), require repeated checks of existing behaviour.\footnote{Angular Material v17 migration guide: \url{https://v17.material.angular.dev/guide/mdc-migration}.}

\section{Adding Playwright to PAWS}

Playwright is an end-to-end testing framework that controls a browser and checks an application through its rendered user interface. Tests perform actions such as navigation, form input and file uploads, then use assertions to check expected results and detect regressions: previously working behaviour broken by application changes. Playwright supports Chromium, Firefox and WebKit; this PAWS suite is configured for Chromium.

Locators identify elements in the rendered document object model (DOM). The tests use user-facing locators where practical; otherwise, \texttt{data-testid} attributes provide dedicated identifiers in the rendered HTML when visible text or component structure is ambiguous. Before clicking, Playwright waits for the element to be visible, stable, enabled and able to receive events. Retrying assertions wait for expected results up to a timeout.\footnote{Playwright documentation: \url{https://playwright.dev/docs/locators}, \url{https://playwright.dev/docs/actionability} and \url{https://playwright.dev/docs/browsers}.}

In this example, the first \texttt{await} waits for the click operation to complete; the second waits for the dialog's visibility assertion to pass, failing if its timeout expires:
\begin{verbatim}
await page.getByTestId("add-observation-band-button").click();
await expect(
  page.getByTestId("add-observation-band-dialog")
).toBeVisible();
\end{verbatim}

In PAWS, a batch groups proposals for a particular observing call; the tests use a configured test batch. Reusable helpers set up data-policy acceptance in the test browser, select this batch, create or clone proposals, update proposal state and remove test data. Figure~\ref{fig:1} shows how browser actions exercise the Angular interface and API helpers perform selected setup, verification and cleanup operations through the development proxy. The backend and its dependencies must be available during execution.

\articlefigure[width=0.85\textwidth]{C502_f1.png}{fig:1}{Local PAWS test architecture. Browser actions and API helpers reach the backend through the development proxy. Browser isolation does not remove stored proposal data or uploaded files.}

\section{Workflow and Regression Coverage}

Table~\ref{tab:coverage} summarises the submission checks exercised across multiple tests. The tested target path allows authors to proceed without specifying individual observing targets. Tests reach the final submission step without clicking Submit.

\begin{table}[t]
\caption{Main submission workflows tested with Playwright.}
\label{tab:coverage}
\centering
\small
\begin{tabular}{@{}p{0.25\textwidth}p{0.70\textwidth}@{}}
\hline
\textbf{Workflow area} & \textbf{What was tested} \\
\hline
Home/navigation & Opening PAWS and navigating to the submission list \\
Proposal creation & Opening the creation dialog and creating a test-batch draft \\
Team details & Entering principal investigator and technical lead details, including the same-as-PI option \\
Science case & Selecting categories, entering an abstract, declaring AI use and uploading a science case PDF \\
Targets & Proceeding without individual observing targets and checking the guidance message \\
Observation setup & Entering observation details and selecting a band \\
Data management & Selecting a facility, transfer mode and pipeline, and uploading a data management PDF \\
Final submission & Entering additional information without clicking Submit \\
\hline
\end{tabular}
\end{table}

The checks exercise the proposal information described in earlier PAWS work \citep{SchwartzBaloyiSithole2024PAWSSubmission}, alongside validation of required fields, email addresses, title length and PDF uploads. Reusable test PDFs, known as fixtures, check document acceptance, page limits and link warnings: a selected Google Drive link triggers a warning, while a safe-link example does not.

Separate tests cover cloning and revision. Cloning lets authors reuse information from an earlier proposal, including team details, to avoid retyping it. The current test verifies new identifiers, the target batch, draft state, a newly entered title and retention of the source abstract; it does not explicitly check copied team details. Revision allows a submitted proposal to be reopened for editing during the submission period \citep{SchwartzBaloyiSithole2024PAWSSubmission}. The revision test prepares a submitted proposal through the API, reopens it through the interface, edits and saves its title, and verifies the resulting backend state.

The wider suite includes Panel Review tests for proposal lists, review information, report fields, downloads and navigation. This paper focuses on submission workflows.


\section{Test Data, Cleanup and Reliability}

Tests create proposals explicitly marked as test data in a configured test batch against the real backend. Cleanup first marks each proposal as deleted without removing its stored data, then uses a dedicated API endpoint to remove its record from Elasticsearch, the database used to store proposal data, and its associated files. This endpoint is restricted to development and test environments and requires the proposal to be marked, soft-deleted and owned by the current submitter. Failed or interrupted tests can still leave data behind.

Separate browser contexts isolate cookies and local storage, but do not remove stored proposals or isolate the shared backend. Tests modifying the same proposal, user or batch can therefore interfere with one another. Reliable execution requires explicit test-data management as well as browser isolation.


\section{Lessons, Benefits and Limitations}

The main challenge was making tests reliable across asynchronous saves, dialogs, uploads and backend state changes. Tests must wait for observable outcomes rather than assume that a button click completes an operation. Stable selectors and reusable fixtures supported maintainability. The suite is configured to capture screenshots and retain execution traces when tests fail, supporting failure diagnosis. Implementing these workflows also clarified the connections between the interface, API and stored proposal data.

The suite provides executable examples of selected submission workflows and repeatable checks of application changes. Coverage remains incomplete: end-to-end tests are slower than unit tests and depend on a working frontend, backend and prepared test data. Failures may reflect environment or data problems rather than application regressions. The suite complements other testing levels and human review; its effect on review time and regressions detected has not yet been measured.


\section{Next Steps}

Future work should map user acceptance testing (UAT) requirements to automated checks, revealing gaps across inputs, roles and proposal states. Coverage should be measured by features rather than inferred from test counts. Priorities include target and observation validation, proposal constraints, permissions, and reviewer workflows for editing, saving and submission \citep{SchwartzBaloyi2024PAWSReview}.

Further work should improve cleanup after interrupted runs and make automated execution reproducible through consistent backend setup, test identities and prepared proposal data. Evaluation should track regressions detected before release, manual checking time, intermittent test failures and maintenance effort.

Adding nightly automated execution of these tests to the Continuous Integration / Continuous Deployment (CI/CD) pipeline would also be of great benefit.


\section{Conclusion}

This work builds confidence in selected PAWS proposal-author workflows through repeatable browser-level regression checks. Reusable helpers, fixtures and explicit cleanup support maintainability as the application evolves. Passing tests provide evidence for the scenarios exercised, not a guarantee of correctness. The suite complements human review and other testing levels, providing a foundation for broader coverage.

In today's fast-paced software development environment, where AI-powered tools are accelerating delivery cycles, automated regression testing is more critical than ever to ensure that existing user-facing features remain fully functional with each new deployment.

\bibliography{C502}

\end{document}
